% Кейс 3. Генерация продуктовых гипотез
% Компиляция: XeLaTeX или LuaLaTeX (рекомендуется), либо pdfLaTeX. Графики берутся из папки fig/ рядом с файлом.
\documentclass[11pt,a4paper]{article}

\usepackage{iftex}
\ifPDFTeX
  \usepackage[T2A]{fontenc}
  \usepackage[utf8]{inputenc}
  \usepackage[english,russian]{babel}
\else
  \usepackage{fontspec}
  \usepackage{polyglossia}
  \setdefaultlanguage{russian}
  \setotherlanguage{english}
  \IfFontExistsTF{CMU Serif}{\setmainfont{CMU Serif}\setsansfont{CMU Sans Serif}\setmonofont{CMU Typewriter Text}}{%
  \IfFontExistsTF{DejaVu Serif}{\setmainfont{DejaVu Serif}\setsansfont{DejaVu Sans}\setmonofont{DejaVu Sans Mono}}{%
  \IfFontExistsTF{Liberation Serif}{\setmainfont{Liberation Serif}\setsansfont{Liberation Sans}\setmonofont{Liberation Mono}}{%
  \IfFontExistsTF{Times New Roman}{\setmainfont{Times New Roman}\setsansfont{Arial}\setmonofont{Courier New}}{%
  \setmainfont{Noto Serif}\setsansfont{Noto Sans}\setmonofont{Noto Sans Mono}}}}}
\fi

\usepackage[margin=2.2cm]{geometry}
\usepackage{graphicx}
\usepackage{booktabs}
\usepackage{tabularx}
\usepackage{array}
\usepackage{float}
\usepackage{caption}
\usepackage{enumitem}
\usepackage[table]{xcolor}
\usepackage{hyperref}
\hypersetup{colorlinks=true, linkcolor=black, urlcolor=blue!50!black}
\setlist{itemsep=2pt, topsep=4pt}
\captionsetup{font=small, labelsep=period}
\renewcommand{\arraystretch}{1.15}
\setlength{\parskip}{3pt}
\setlength{\emergencystretch}{2.5em}

% Рисунок: если файла нет, вместо него рамка с подсказкой (документ собирается всегда)
\newcommand{\fig}[3]{%
  \begin{figure}[H]\centering
  \IfFileExists{fig/#1.png}{\includegraphics[width=#2\linewidth]{fig/#1.png}}%
    {\fbox{\parbox[c][3cm][c]{#2\linewidth}{\centering\small Нет файла \texttt{fig/\detokenize{#1}.png}.\\Скопируйте его из папки \texttt{case3/fig} на Google Drive.}}}
  \caption{#3}\label{fig:#1}
  \end{figure}}

% Блоки карточки гипотезы
\newcommand{\hblock}[1]{\par\smallskip\noindent\textbf{#1.}~}

\title{Кейс 3. Генерация продуктовых гипотез\\[2pt]\large подписочный VPN-сервис: дамп серверной БД и логи инцидентов}
\author{Производственная практика}
\date{}

\begin{document}
\maketitle
\tableofcontents
\newpage

%==============================================================
\section{Задача и подход}

Владелец интернет-сервиса (клиент-серверное мобильное приложение для Android, iOS и Huawei, продаётся подписками по тарифам) передал дамп серверной БД, выгрузку логов клиентов по инцидентам и структуру таблиц. Требуется сопоставить данные о клиентах и их действиях, сформулировать продуктовые гипотезы и подтвердить или опровергнуть каждую на этих данных. Каждая гипотеза оценивается отдельно, поэтому приоритет отдан проверяемости, а не числу гипотез.

Проверено 10 гипотез: четыре из условия (H1--H4) и шесть собственных (H5--H10). Каждая оформлена одинаково: утверждение, данные, метод, результат, ограничения, вывод, рекомендация. Сводка вердиктов --- в табл.~\ref{tab:summary}.

\paragraph{Инструменты и воспроизводимость.} Python в Google Colab: DuckDB для запросов к таблицам объёмом до 22,7~млн строк (поверх Parquet), pandas и SciPy для статистики, matplotlib для графиков. Решение разбито на пять ноутбуков, каждый запускается целиком и скачивает данные по публичной ссылке владельца:
\begin{itemize}
  \item \texttt{case3\_00\_load\_recon} --- загрузка CSV в Parquet с обезличиванием, первичная разведка;
  \item \texttt{case3\_01\_recon\_extra} --- часовые пояса и проверки для витрины;
  \item \texttt{case3\_02\_mart\_h3} --- витрина «один пользователь --- одна строка» и гипотеза H3;
  \item \texttt{case3\_03\_db\_hypotheses} --- гипотезы по БД: H1, H2 (реклама), H5--H9;
  \item \texttt{case3\_04\_logs\_hypotheses} --- клиентские логи: H2 (трафик), H4, H10, а также уточнения по БД.
\end{itemize}

\paragraph{Статистика.} Доли приводятся с 95\,\% доверительным интервалом Уилсона, сравнение двух долей --- двусторонним $z$-тестом, сравнение долей между версиями приложения --- критерием $\chi^2$. Время до первой покупки оценивается кривой Каплана--Мейера, потому что у недавно зарегистрированных пользователей покупка ещё могла не наступить (цензурирование). Все данные наблюдательные: выводы сформулированы как связь, а не как причина, если не оговорено иное.

\paragraph{Роль ИИ-ассистента.} Код ноутбуков, план проверки гипотез и черновик интерпретации подготовлены с помощью ИИ-ассистента Claude (Anthropic). Ассистент не имел доступа к самим данным: ноутбуки запускались автором в Colab, ассистенту передавались только агрегированные результаты (счётчики, доли, квантили). Автор определял состав и объём работы, запускал вычисления, проверял результаты и принимал решения по спорным местам (переформулировка H2, набор гипотез, формат сдачи). Ошибки, найденные по ходу (неверные алиасы SQL, неправильный признак автопродления в первой версии H7, учёт пустых значений в долях), исправлены и пересчитаны.

%==============================================================
\section{Данные и их подготовка}

\subsection{Состав и объём}
\begin{itemize}
  \item 1\,844\,880 пользователей (с 25.11.2024), из них 1\,176\,913 гостей (гостевой режим с июня 2026) и 667\,962 зарегистрированных клиента;
  \item 136\,006 заказов, 124\,863 транзакции, 69\,968 подписок; провайдеры TipTopPay, Yookassa (основной), Heleket;
  \item 22,75~млн записей о подключениях к VPN-нодам за 01.02--29.09.2026 (месячные файлы, март и апрель в одном); 87 нод в 252 локациях;
  \item 7,55~млн просмотров рекламы, 3,05~млн сессий, 2,38~млн устройств;
  \item клиентские логи: 926 архивов (инцидентов) от 787 пользователей, 62\,550 файлов, 44~ГБ в распакованном виде: журнал приложения, журналы VPN-ядра (Xray), системный журнал Android.
\end{itemize}

\subsection{Особенности выгрузки и как они учтены}
\begin{enumerate}
  \item \textbf{В CSV нет строки заголовков} (кроме логов подключений). Имена колонок взяты из \texttt{db\_structure}; разделитель задан явно (\texttt{;}), иначе JSON-поля с запятыми ломают разбор.
  \item \textbf{Разные часовые пояса.} Поля без пояса в \texttt{users}, логах подключений, рекламе и прокси-логах записаны по московскому времени, в \texttt{orders}, \texttt{transactions}, \texttt{subscribe} --- в UTC. Определено по совпадению суточного профиля активности с эталоном (\texttt{sessions}, поле с поясом); без поправки 65\,\% первых заказов оказываются «раньше регистрации», с поправкой --- 0,01\,\%. Всё приведено к UTC.
  \item \textbf{Персональные данные.} При загрузке удалены пароли, имена, push-токены, данные карт, ссылки на оплату; IP, идентификаторы устройств и внешних аккаунтов заменены хешем с секретной солью; от карты оставлен только отпечаток (хеш BIN, последних 4 цифр и срока). В отчёте только агрегаты.
  \item \textbf{Нет объёма трафика в БД и таблицы лимитов} (\texttt{sub\_limit} отсутствует в выгрузке). Потребление трафика восстановлено по клиентским логам (H2).
  \item \textbf{Целостность ключей} высокая: внешние ключи находятся в родительских таблицах в 99,99--100\,\% случаев; номер папки логов совпадает с \texttt{users.id} у всех 787 пользователей.
\end{enumerate}

\subsection{Определения}
\begin{itemize}
  \item \textbf{Бесплатное использование} --- подключение без \texttt{subscribe\_id}. Проверка: в момент таких подключений активная платная подписка есть у 0,0--0,2\,\% пользователей; все они идут через ноды бесплатного пула.
  \item \textbf{Триал} --- платный: 10--30\,₽ за 3 дня тарифа Lite (\texttt{orders.trial}); момент оплаты совпадает с \texttt{users.trial\_used\_at} у 22\,192 из 22\,193 пользователей.
  \item \textbf{Покупка} --- оплаченный обычный (не триальный) заказ; \textbf{первая покупка} --- первый такой заказ пользователя.
  \item \textbf{Автосписание} --- \texttt{orders.pre\_payment}; этот флаг стоит только у продлений. \textbf{Автопродление подписки} --- \texttt{subscribe.pre\_payment}.
  \item Платить может только зарегистрированный клиент: у гостей нет ни одного заказа.
\end{itemize}

\subsection{Общие ограничения}
\begin{itemize}
  \item 27.02.2026 число подключений упало на 98\,\% одновременно на бесплатных и платных нодах и восстанавливалось несколько недель: это общий сбой сервиса или выгрузки. Сравнения «месяц к месяцу» по подключениям поэтому не делаются.
  \item С июня 2026 большинство новых пользователей --- гости, и состав «клиентов» изменился: в клиенты теперь чаще попадают те, кто регистрируется ради оплаты. Сравнения когорт до и после июня оговариваются отдельно.
  \item Клиентские логи прислали только пользователи с проблемами; 94\,\% этих пользователей --- на Android; архивы собраны 13--28.09.2026, файлы внутри в основном за август--сентябрь 2026. Частоты по логам описывают инциденты, а не всю аудиторию.
  \item Страна пользователя в данных отсутствует; отключение бесплатного тарифа для RU, KZ и US с 19.03.2026 (\texttt{country\_disabled\_products}) проверить по пользователям нельзя, бесплатное использование после этой даты продолжается.
\end{itemize}

%==============================================================
\section{Сводка гипотез}

\begin{table}[H]
\centering\small
\caption{Вердикты по гипотезам}\label{tab:summary}
\begin{tabularx}{\linewidth}{@{}l>{\raggedright\arraybackslash}X>{\raggedright\arraybackslash}p{3.0cm}>{\raggedright\arraybackslash}X@{}}
\toprule
& Гипотеза & Вердикт & Ключевой результат \\
\midrule
H1 & Мультиаккаунты ради триала, и они вредят конверсии & опровергнута в сильной форме & повторные триалы: 3,8\,\% всех триалов (нижняя оценка); конверсия как у одиночных (55,4 vs 55,3\,\%, $p=0{,}97$) \\
H2 & Лимит 500~МБ/день поставлен «на глаз»; какой объём оптимален & частично подтверждена & медиана дня 180~МБ, лимит упирается примерно в 30\,\% дней; один просмотр рекламы даёт +200~МБ \\
H3 & Покупка всегда после бесплатного периода и триала & опровергнута & 36\,\% первых покупок без триала; 77\,\% «покупок после триала» --- автосписание в конце триала \\
H4 & Есть поведенческие паттерны, ведущие к типовым ошибкам & подтверждена & 74\,\% инцидентов: API-домен не разрешается DNS при старте приложения $\Rightarrow$ сбой конфигурации и WebSocket \\
H5 & Неуспешные платежи --- крупная потеря выручки & подтверждена & 46\,\% первых оплат неуспешны, 88\,\% отказов --- нехватка средств; сбой продлений в июне: 2,47~млн\,₽ \\
H6 & Длинные периоды не продаются, несмотря на скидки & подтверждена & 0,9\,\% заказов; 74--94\,\% первых попыток неуспешны \\
H7 & Автопродление удерживает подписчиков & подтверждена (связь) & продление 72\,\% против 27\,\% без автопродления \\
H8 & Гости дают охват, но монетизация --- через регистрацию & подтверждена & в клиенты переходят 2,6--3,0\,\% гостей; бывшие гости платят в 3 раза чаще прямых клиентов \\
H9 & Бесплатный пул нод перегружен и концентрирован & подтверждена, часть про обвал 27.02 опровергнута & до 26 тыс. пользователей на бесплатную ноду в день против 77 на платной \\
H10 & Старые версии приложения дают больше ошибок & слабо подтверждена & недоступность API: 95 и 86\,\% в версиях 2026.7.16 и 2026.8.19 против 74\,\% в 2026.9.9 ($p=0{,}005$) \\
\bottomrule
\end{tabularx}
\end{table}

%==============================================================
\section{Гипотезы из условия}

%--------------------------------------------------------------
\subsection{H1. Мультиаккаунты ради триала}

\hblock{Утверждение} Часть пользователей заводит новые аккаунты на том же устройстве или с той же картой, чтобы снова взять триал; такие повторные триалы реже превращаются в покупку.

\hblock{Данные} \texttt{devices.device\_identifier} (заполнен у 42\,\% устройств), отпечатки карт из \texttt{saved\_payment\_methods}, \texttt{orders} (триалы и покупки), \texttt{sessions.ip\_address} (хеш), витрина пользователей.

\hblock{Метод} Аккаунты объединены в группы, если у них общий идентификатор устройства или общий отпечаток карты (группы больше 50 аккаунтов отброшены как технические идентификаторы). Внутри группы оплаченные триалы упорядочены по времени: первый --- «первичный», остальные --- «повторные». Конверсия триала в покупку за 30 дней сравнивается для повторных, первичных и одиночных триалов (только триалы старше 30 дней). Общий IP проверен отдельно как слабый признак: учитывались только IP, за которыми не больше 5 аккаунтов, чтобы исключить мобильные сети с общим адресом.

\hblock{Результат}
\begin{itemize}
  \item Групп по устройству --- 8\,551, по карте --- 322; всего 26\,464 аккаунта в 7\,490 группах (97\,\% связей --- через устройство).
  \item Из 22\,193 оплаченных триалов 852 --- повторные (\textbf{3,8\,\%}), в 486 группах; одна группа содержит 79 триалов (вероятно, эмулятор или общее устройство). Медиана между соседними триалами в группе --- 15,9 дня.
  \item Конверсия за 30 дней (рис.~\ref{fig:h1_conversion}): одиночные --- 55,3\,\% (ДИ 54,6--56,0, $n=17\,574$), повторные --- 55,4\,\% (51,6--59,1, $n=681$), первичные в группе --- 60,8\,\% (56,0--65,3, $n=423$). Разница повторных и одиночных $+0{,}07$~п.п., $p=0{,}97$.
  \item Косвенные признаки «аккаунта под триал» у повторных не сильнее: триал в первый час после регистрации --- 48,9\,\% против 57,6\,\% у одиночных; активность только в окне триала без покупки --- 13,4 против 13,6\,\%.
  \item Доля повторных выше на iOS (12,6\,\%), Windows (11,7\,\%) и Android TV (11,0\,\%), чем на Android (2,6\,\%) --- вероятно, сказывается разное поведение идентификатора устройства на платформах.
  \item Слабый признак: 2\,569 IP с двумя и более триальщиками охватывают 5\,412 триальщиков (24\,\%); сюда попадают семьи и общий домашний Wi-Fi.
\end{itemize}

\fig{h1_conversion}{0.8}{H1. Конверсия триала в покупку за 30 дней с 95\,\% ДИ}

\hblock{Ограничения} Идентификатор устройства есть лишь у 42\,\% устройств, поэтому 3,8\,\% --- нижняя оценка; отпечаток карты есть только у 5,8 тыс. из 22 тыс. триальщиков (остальные, по-видимому, платят не картой). IP-признак шумный и в оценку не включён.

\hblock{Вывод} Злоупотребление триалом существует, но мало (порядка 4\,\% триалов по надёжным признакам), а повторные триалы конвертируются в покупку так же, как обычные. Платный триал (10--30\,₽ с привязкой способа оплаты) уже работает как барьер. Гипотеза о вредном массовом мультиаккаунтинге \textbf{опровергается}.

\hblock{Рекомендация} Не ужесточать выдачу триала: выигрыш мал, а риск отсечь платящих есть. Достаточно мониторинга: число триалов на идентификатор устройства и отпечаток карты с порогом, например, не больше двух триалов на устройство в квартал.

%--------------------------------------------------------------
\subsection{H2. Лимиты бесплатного тарифа}

\hblock{Утверждение} Лимит 500~МБ в день (15~ГБ в месяц) выбран «на глаз». Исходная формулировка («подобрать оптимальные объёмы по реальному потреблению») переформулирована, так как объёма трафика в БД нет: \emph{какая доля дней бесплатного использования упирается в лимит, сколько потребляют за день и как используется расширение лимита рекламой}.

\hblock{Данные} Клиентские логи: строки \texttt{[TrafficLimitBloc] Traffic info changed: \ldots{} bytes limit, \ldots{} bytes used} и \texttt{limitReached}; \texttt{adv\_view\_history}, логи подключений без подписки, \texttt{mobile\_version} (флаг \texttt{showRewardedAdLimit} --- реклама за лимит).

\hblock{Метод} По логам: для каждого пользователя и дня --- максимальное «использовано» и лимит в начале и в конце дня; квантили потребления; доля дней, где потребление достигло лимита; рост лимита за день; «что было бы» при лимите $L$ --- доля дней, где наблюдаемое потребление не меньше $L$. По БД: доля дней бесплатного использования с просмотром рекламы и связь рекламы с покупкой. Чтобы покупка не влияла на просмотры, взяты клиенты, не купившие в первые 7 дней: просмотры считаются в дни 1--7, покупка --- в дни 8--37.

\hblock{Результат}
\begin{itemize}
  \item Счётчики есть у 80 пользователей (405 пользователе-дней). Базовый лимит --- 500~МиБ (у 72 из 80); лимит растёт шагами ровно по 200~МБ: 500 $\to$ 691 $\to$ 881 $\to$ 1072~МиБ, то есть \textbf{один просмотр рекламы добавляет 200~МБ}.
  \item Потребление за день (рис.~\ref{fig:h2_daily_usage}): медиана \textbf{180~МБ}, квартили 22 и 624~МБ, 90-й процентиль 1,38~ГБ, 95-й --- 2,07~ГБ.
  \item Потребление достигло базового лимита в \textbf{28\,\%} дней, 80\,\% лимита --- в 30\,\%; лимит расширялся рекламой в 37\,\% дней, медиана прироста 777~МБ за день (около 4 просмотров).
  \item Доля дней, где потребление не меньше $L$: 300~МБ --- 38,8\,\%; 500~МБ --- 30,4\,\%; 750~МБ --- 23,0\,\%; 1~ГБ --- 17,3\,\%; 2~ГБ --- 5,4\,\%.
  \item По всей базе (БД) реклама есть в 18--33\,\% дней бесплатного использования, и доля снижается: февраль 33\,\%, сентябрь 18\,\% (рис.~\ref{fig:h2_ads_share}); 69\,\% бесплатных пользователей не смотрели рекламу ни разу.
  \item Связь с покупкой положительная: покупка в дни 8--37 при 0 просмотрах --- 0,46\,\% (ДИ 0,43--0,50), при 1--5 --- 0,64\,\%, при 6--20 --- 0,91\,\%, при 20+ --- 1,12\,\% (0,91--1,38).
\end{itemize}

\fig{h2_daily_usage}{0.8}{H2. Суточное потребление на бесплатном тарифе по клиентским логам}
\fig{h2_ads_share}{0.8}{H2. Доля дней бесплатного использования с просмотром рекламы}

\hblock{Ограничения} Выборка мала и смещена: 80 пользователей, приславших логи из-за проблем, почти все на Android, данные в основном за август--сентябрь 2026. Потребление выше текущего лимита не наблюдается, поэтому доли для $L>500$~МБ --- нижние границы. Счётчик сбрасывается не по локальной дате устройства (сброс виден лишь в 49\,\% переходов между датами), поэтому «потребление за день» --- приближение. Месячный лимит 15~ГБ оценить не удалось: строк с месячным лимитом почти нет.

\hblock{Вывод} Лимит 500~МБ \textbf{не завышен и не занижен «на глаз»}: он втрое выше медианного дня и не мешает примерно 70\,\% дней, но ограничивает 30\,\% дней, и примерно в этих днях пользователи расширяют его рекламой по 200~МБ за просмотр. Реклама не заменяет покупку: активные зрители покупают чаще (это признак вовлечённости, а не доказанное следствие рекламы). Однозначного «оптимального» значения данные не дают: это выбор между доходом от рекламы (ниже лимит --- больше дней с рекламой) и комфортом пользователя. Гипотеза \textbf{частично подтверждена}.

\hblock{Рекомендация} Проверить A/B-тестом лимиты 300 / 500 / 750~МБ по метрикам: просмотры рекламы на активного пользователя, удержание через 7 и 30 дней, переход в клиенты и покупка. Ориентир по логам: снижение до 300~МБ увеличит долю дней с упором в лимит с 30 до примерно 39\,\%, повышение до 750~МБ уменьшит до 23\,\%.

%--------------------------------------------------------------
\subsection{H3. Путь от бесплатного тарифа к покупке и роль триала}

\hblock{Утверждение} Пользователи сначала пользуются бесплатным тарифом, затем берут триал и только потом покупают подписку. Вопросы условия: как долго длится бесплатный период, всегда ли покупке предшествует триал, сколько клиентов покупает без триала.

\hblock{Данные} Витрина пользователей: регистрация, первый оплаченный триал, первая покупка, платформа; \texttt{subscribe} (конец триала); логи подключений и реклама (с 01.02.2026).

\hblock{Метод} Клиенты, зарегистрированные с 10.06.2025 (появление триала) по 28.09.2026, разделены по пути к первой покупке. Конверсия триала --- за 30 дней, только по триалам старше 30 дней. Момент покупки после триала сравнивается с концом триальной подписки с точностью до минут. Доля купивших к дню $t$ --- оценка Каплана--Мейера. Бесплатные дни до покупки --- число дней с подключениями без подписки до первой покупки.

\hblock{Результат}
\begin{itemize}
  \item Из 657\,557 клиентов купили подписку 19\,178 (2,9\,\%): \textbf{через триал --- 12\,271 (64,0\,\%), без триала --- 6\,907 (36,0\,\%, ДИ 35,3--36,7)} (рис.~\ref{fig:h3_paths}). Ещё 9\,541 оплатили только триал.
  \item Конверсия триала в покупку за 30 дней --- 55,2\,\% (ДИ 54,5--56,0); iOS и macOS --- 62\,\%, Android --- 55\,\%, Windows --- 51\,\%.
  \item \textbf{77,4\,\% покупок после триала совершены в пределах $\pm5$ минут от конца триала} (рис.~\ref{fig:h3_trial_to_paid}) --- это автоматическое списание, а не самостоятельное решение пользователя; позже суток после конца триала --- 10,5\,\%.
  \item Покупают быстро: к 30-му дню после регистрации купили 2,5\,\% клиентов, к 180-му --- 3,0\,\% (рис.~\ref{fig:h3_km}). 70\,\% покупок без триала --- в первые сутки после регистрации; 62\,\% купивших не имели ни одного дня бесплатного использования до покупки, 90\,\% --- не больше 6 дней; рекламу до покупки смотрели 17,6\,\%.
  \item С мая 2026 доля покупок без триала выросла с 10--20\,\% до 47--53\,\%: в мае зарегистрировалось 112,7~тыс. клиентов при низкой доле триалов (2,2\,\%), оплата в первый час после регистрации выросла с 20 до 51--77\,\%, с июля это в основном бывшие гости.
  \item Платформа: к 30-му дню купили 15,6\,\% клиентов macOS против 2,3\,\% на Android; на десктопах реклама за лимит выключена, т.\,е. бесплатный режим там беднее (связь, а не доказанная причина).
  \item 9\,388 клиентов пытались оплатить триал, но оплата не прошла ни разу.
\end{itemize}

\fig{h3_paths}{0.8}{H3. Путь к первой покупке}
\fig{h3_trial_to_paid}{0.8}{H3. Время от оплаты триала до первой покупки}
\fig{h3_km}{0.8}{H3. Доля клиентов, купивших подписку, по кварталу регистрации (Каплан--Мейер)}

\hblock{Ограничения} С июня 2026 гости не входят в число клиентов, поэтому рост конверсии свежих когорт (5,1\,\% к 30-му дню в III квартале 2026 против 1,2\,\% в I квартале) в значительной мере объясняется сменой состава клиентов, а не улучшением продукта. Изменение онбординга в мае 2026 выведено косвенно (объёмы и время оплаты), прямых данных о релизе нет.

\hblock{Вывод} Гипотеза «покупка всегда после триала» \textbf{опровергается}: 36\,\% первых покупок --- без триала, а с мая 2026 --- около половины. Длительного бесплатного периода перед покупкой обычно нет: медиана --- 0 бесплатных дней. Триал работает как «пробный период с привязкой оплаты и автосписанием», а не как этап знакомства с продуктом.

\hblock{Рекомендация} Считать в отчётности две метрики раздельно: автосписание после триала и самостоятельную покупку. Сосредоточить конверсионные механики на первых сутках после регистрации. Для пользователей с неуспешной оплатой триала (9,4~тыс.) предложить альтернативный способ оплаты.

%--------------------------------------------------------------
\subsection{H4. Паттерны поведения, приводящие к типовым ошибкам}

\hblock{Утверждение} В логах инцидентов есть устойчивые сценарии, после которых возникают типовые ошибки, и большинство инцидентов связано не с VPN-нодами, а с недоступностью сервера приложения.

\hblock{Данные} 926 архивов (инцидентов) от 787 пользователей: журнал приложения (Flutter), журналы VPN-ядра Xray, системный журнал Android; роль и платежи пользователя из БД.

\hblock{Метод} Разбор всех архивов регулярными выражениями; для каждого инцидента --- наличие ошибок 11 типов, их сочетания, а также три события перед первой ошибкой в каждом файле журнала приложения (цифры в шаблонах заменены).

\hblock{Результат}
\begin{itemize}
  \item Почти в каждом инциденте есть сетевые исключения (97\,\%), тайм-ауты (95\,\%) и ошибки HTTP-клиента (92\,\%) --- это фон, по которому инциденты не различаются.
  \item Различающий признак: \textbf{домен API (\texttt{api.pipster.win}) не разрешается DNS в 74,3\,\% инцидентов}; сбои WebSocket --- 79\,\%; сбои DoH-запросов VPN-ядра (dns.google, cloudflare, quad9) --- 79\,\% (рис.~\ref{fig:h4_error_types}).
  \item Сбой DNS для API и сбой DoH в VPN-ядре идут вместе: при недоступном API DoH-сбои есть в 86,8\,\% инцидентов, при доступном --- в 56,7\,\%.
  \item Самое частое сочетание (287 инцидентов, 31\,\%) --- полный набор: API не разрешается, WebSocket, тайм-ауты, DoH, тайм-ауты и отказы соединений VPN-ядра.
  \item Сценарии перед первой ошибкой:
  \begin{enumerate}
    \item подключение WebSocket к API: заголовки $\to$ ожидание $\to$ \emph{Handshake failed} (около 8,9~тыс. файлов);
    \item запуск приложения: инициализация VPN-адаптера и SDK $\to$ \emph{загрузка конфигурации (S3) не удалась} (около 5,0~тыс.);
    \item VPN успешно включён $\to$ \emph{не удалось запланировать уведомление} (около 1,9~тыс.);
    \item смена сети $\to$ исключение при запуске VPN или тайм-аут при остановке (около 250).
  \end{enumerate}
  \item На Android 13--16 недоступность API встречается в 74--86\,\% инцидентов, на Android 10 --- в 47\,\%.
  \item Гости и клиенты не различаются по профилю ошибок.
\end{itemize}

\fig{h4_error_types}{0.85}{H4. Доля инцидентов с ошибкой данного типа}

\hblock{Ограничения} Логи присылают только пользователи с проблемами; архивы собраны во второй половине сентября 2026, пользователи почти все на Android. Причина недоступности DNS (блокировка, провайдер, настройки «частного DNS» на устройстве) по логам не устанавливается.

\hblock{Вывод} Гипотеза \textbf{подтверждается}. Типовой сценарий инцидента: пользователь запускает приложение или меняет сеть при выключенном VPN, домен API и публичные DoH-резолверы недоступны, приложение не получает конфигурацию и не может открыть WebSocket, а значит, не может подключиться. Ноды при этом не главная причина.

\hblock{Рекомендация} Резервный путь получения конфигурации при недоступном API: кэш последней конфигурации, встроенный список запасных адресов, резервные домены и DoH-резолверы. Отдельно --- исправить планирование уведомлений после включения VPN (вероятно, разрешение на уведомления в новых версиях Android).

%==============================================================
\section{Собственные гипотезы}

%--------------------------------------------------------------
\subsection{H5. Неуспешные платежи}

\hblock{Утверждение} Значительная часть оплат не проходит, главным образом из-за нехватки средств, и часть пользователей после неудачи не возвращается.

\hblock{Данные} \texttt{orders} (тип заказа: триал, первая покупка, продление), \texttt{transactions} (причины отказа Yookassa).

\hblock{Метод} Доля неуспешных заказов по типу, сумме и периоду; причины отказа; что происходит после неуспешного заказа (оплата в течение часа, 7 дней, никогда --- только заказы старше 7 дней); недельная динамика зависших (\texttt{pending}) заказов.

\hblock{Результат}
\begin{itemize}
  \item Неуспешны \textbf{46,2\,\% первых покупок}, 34,3\,\% продлений, 41,3\,\% триалов. Доля растёт с суммой первой покупки: до 200\,₽ --- 42,6\,\%, 300--499\,₽ --- 58,4\,\%, 500--999\,₽ --- 62,3\,\%, от 1000\,₽ --- 92,7\,\% (рис.~\ref{fig:h5_failed_by_period}).
  \item Причина отказа: \textbf{нехватка средств --- 87,9\,\%} (первые покупки) и 80,0\,\% (продления); отзыв разрешения на списание --- 7,7 и 10,1\,\%.
  \item После неуспешной первой покупки в течение часа платят 16,9\,\%, в течение 7 дней --- 30,7\,\%, \textbf{никогда --- 64,1\,\%}; после неуспешного триала никогда не платят 69,0\,\%. Первую покупку так и не смогли оплатить 9\,482 из 27\,997 пытавшихся (33,9\,\%); верхняя оценка упущенного первого платежа --- 5,6~млн\,₽.
  \item \textbf{Сбой продлений 25.05--28.06.2026}: 60--77\,\% продлений еженедельно оставались в статусе \texttt{pending} при обычных 0--1\,\%. Зависло 8\,353 продления 4\,097 пользователей на 2,47~млн\,₽; в течение 14 дней оплатили 45,2\,\%, \textbf{никогда не оплатили 48,0\,\%}. Для сравнения: среди всех продлений апреля 2026 больше не платили около 36\,\%.
\end{itemize}

\fig{h5_failed_by_period}{0.8}{H5--H6. Доля неуспешных заказов по периоду подписки}

\hblock{Ограничения} Причина отказа есть только у Yookassa и только при отмене банком; у других провайдеров и при уходе со страницы оплаты причина неизвестна. Оценка потерь от сбоя --- сравнение с обычным месяцем, а не контролируемый эксперимент.

\hblock{Вывод} Гипотеза \textbf{подтверждается}: платёжеспособность --- главный барьер, и неуспешная оплата чаще всего заканчивает путь пользователя. Июньский сбой продлений, по грубой оценке, увеличил долю ушедших примерно на 12~п.п., то есть около 500 подписчиков сверх обычного.

\hblock{Рекомендация} Повторные попытки списания продлений с учётом дня (ближе к датам зарплат), предложение альтернативного способа оплаты сразу после отказа, более дешёвый короткий вариант для отказов по нехватке средств, мониторинг доли \texttt{pending} с оповещением; кампания возврата для 4,1~тыс. пользователей из сбоя.

%--------------------------------------------------------------
\subsection{H6. Длинные периоды подписки}

\hblock{Утверждение} Скидки 15, 30 и 60\,\% на 6, 12 и 24 месяца почти не работают: длинные периоды продаются редко, а попытки их оплатить чаще проваливаются.

\hblock{Данные} \texttt{orders} (период, статус, сумма), \texttt{product\_period\_discounts}.

\hblock{Метод} Доля периодов в оплаченных заказах и в выручке Yookassa; доля неуспеха по периоду; переход на месячную подписку в течение суток после неуспешной попытки длинной.

\hblock{Результат}
\begin{itemize}
  \item Длинные периоды --- \textbf{0,9\,\% оплаченных заказов} (260 на 6 месяцев, 132 на год, 23 на 2 года), но 8,6\,\% выручки Yookassa (1,22~млн из 14,2~млн\,₽). Максимум доли --- декабрь 2025 (2,6\,\%), с весны 2026 --- 0,5--1,3\,\%.
  \item Неуспешна первая попытка: 6 месяцев --- 74,3\,\%, год --- 81,5\,\%, 2 года --- 93,8\,\% (месяц --- 44,7\,\%); у продлений --- 56,9 / 67,5 / 82,5\,\% против 34,1\,\%.
  \item После неуспешной попытки в течение суток месячную подписку оплачивают 34,6\,\% (6 месяцев), 23,6\,\% (год), 15,3\,\% (2 года); тот же длинный период --- 8,2 / 4,5 / 2,6\,\%.
\end{itemize}

\hblock{Ограничения} У неуспешных длинных заказов почти нет причины отказа банка (2 записи), поэтому нельзя сказать, отказал ли банк или пользователь ушёл со страницы оплаты. Связь с суммой видна косвенно --- по росту доли неуспеха с суммой (H5).

\hblock{Вывод} Гипотеза \textbf{подтверждается}: длинные периоды приносят заметную долю выручки при крошечном числе заказов, но крупная разовая сумма, по-видимому, не проходит, и часть пользователей переходит на месяц.

\hblock{Рекомендация} Промежуточный период на 3 месяца, отображение цены в пересчёте на месяц, предложение длинного периода тем, кто уже продлил подписку два-три раза (у них выше вероятность иметь нужную сумму и желание экономить).

%--------------------------------------------------------------
\subsection{H7. Автопродление и удержание}

\hblock{Утверждение} Подписчики с включённым автопродлением продлевают подписку заметно чаще.

\hblock{Данные} \texttt{subscribe.pre\_payment} (автопродление подписки), \texttt{orders} (последующие оплаты и автосписания).

\hblock{Метод} Первые месячные покупки, у которых конец подписки плюс 7 дней уже наступил. Продление --- следующая оплаченная обычная подписка не позже 7 дней после конца первой. Удержание --- число продлений у первых покупок до 01.04.2026.

\hblock{Результат}
\begin{itemize}
  \item Автопродление включено у 84\,\% первых месячных подписок (11\,920 из 14\,160).
  \item Продление: \textbf{с автопродлением --- 71,7\,\% (ДИ 70,9--72,5), без --- 26,9\,\% (25,1--28,8)}; в группе с автопродлением 54,1\,\% продлений прошли автосписанием.
  \item Автопродление отключили до конца срока у 54,3\,\% подписок, где оно было включено.
  \item Тариф: Silver продлевают чаще Lite (71,8 против 64,0\,\% без разделения по автопродлению). После триала и без него --- 71,0 и 73,3\,\% (с автопродлением).
  \item Удержание: хотя бы одно продление --- 70,9\,\%, два --- 54,1\,\%, три --- 39,5\,\%, пять и более --- 24,9\,\% ($n=5\,524$).
\end{itemize}

\hblock{Ограничения} Сравнение наблюдательное: автопродление оставляют более мотивированные пользователи, поэтому разница 72 против 27\,\% --- верхняя оценка эффекта самого автопродления.

\hblock{Вывод} Гипотеза \textbf{подтверждается как связь}: подписки с автопродлением продлеваются в 2,7 раза чаще. При этом больше половины отключает автопродление в течение первого месяца.

\hblock{Рекомендация} Снижать отключение автопродления: напоминание о списании заранее, возможность приостановить подписку вместо отключения, ответ на отключение (предложение тарифа подешевле). Эффект проверять A/B-тестом.

%--------------------------------------------------------------
\subsection{H8. Гости и переход в клиенты}

\hblock{Утверждение} Гостевой режим (с июня 2026) дал основной прирост аудитории, но гости не могут платить, и монетизация идёт через регистрацию; долгое бесплатное использование к регистрации и покупке ведёт редко.

\hblock{Данные} \texttt{users} (роль, \texttt{is\_was\_guest}), витрина: бесплатные дни, реклама, покупки.

\hblock{Метод} Переход гостей в клиенты по месяцу создания; покупка у бывших гостей и прямых клиентов, зарегистрированных в июне--августе 2026; активность гостей.

\hblock{Результат}
\begin{itemize}
  \item Гостей --- 1,18~млн (64\,\% всех пользователей); в июле--сентябре 2026 гостей создано 201--508~тыс. в месяц. В клиентов переходят \textbf{2,6--3,0\,\%} гостей каждого месяца.
  \item Бывшие гости платят в 3 раза чаще прямых клиентов: 12,5 против 3,9\,\% (триал --- 10,9 против 3,3\,\%).
  \item Чем дольше бывший гость пользовался бесплатно, тем реже он платит: 0 бесплатных дней --- 30,4\,\%, 1--7 --- 5,6\,\%, 8--30 --- 3,5\,\%, больше 30 --- 2,4\,\%.
  \item 24\,\% гостей ни разу не подключились, 38\,\% подключались один день; рекламу смотрят 2\,\% гостей без подключений, 31\,\% активных 2--7 дней и 86\,\% активных больше 30 дней.
\end{itemize}

\hblock{Ограничения} Неизвестно, что означает \texttt{created\_at} бывшего гостя --- создание гостевой записи или регистрацию; поэтому сроки перехода не оцениваются. Сравнение бывших гостей и прямых клиентов подвержено самоотбору.

\hblock{Вывод} Гипотеза \textbf{подтверждается}: гостевой режим расширил охват, а деньги приносят те, кто быстро регистрируется; длительное бесплатное использование гостем почти не превращается в оплату, но даёт рекламные просмотры.

\hblock{Рекомендация} Предлагать регистрацию в моменты наибольшей ценности (упор в лимит, перед просмотром рекламы, после успешного подключения в первый день). Для долгих гостей --- монетизация рекламой, а не попытки продать подписку.

%--------------------------------------------------------------
\subsection{H9. Нагрузка на бесплатный пул нод}

\hblock{Утверждение} Бесплатный трафик сосредоточен на нескольких нодах, и при росте аудитории нагрузка на одну бесплатную ноду росла быстрее, чем расширялся пул; обвал подключений 27.02.2026 касался только бесплатного пула.

\hblock{Данные} Логи подключений, \texttt{nodes.free\_pool}, \texttt{task\_history}.

\hblock{Метод} Среднее число пользователей на активную ноду в день по месяцам для бесплатного и платного пулов; доля крупнейшей бесплатной ноды; подневная динамика вокруг 27.02.2026.

\hblock{Результат}
\begin{itemize}
  \item Через бесплатный пул идёт 88,3\,\% всех подключений; платные подписчики его не используют.
  \item С февраля по апрель 2026 весь бесплатный трафик обслуживала \textbf{одна нода} (Германия): до 22,3~тыс. пользователей в день. В сентябре --- в среднем 3,4 ноды, \textbf{26~тыс. пользователей на ноду в день} (максимум 37,4~тыс.); на платной ноде --- 77, примерно в 340 раз меньше (рис.~\ref{fig:h9_free_pool_load}).
  \item Доля крупнейшей бесплатной ноды снизилась со 100\,\% до 35\,\%, но нагрузка на ноду выросла в 1,4 раза с февраля.
  \item 27.02.2026 упали оба пула: бесплатный --- с 40,3 до 1,0~тыс. подключений в день, платный --- с 2,0~тыс. до 45; к 08.03 бесплатный восстановился до 40\,\%, платный --- до 62\,\%.
\end{itemize}

\fig{h9_free_pool_load}{0.8}{H9. Пользователи на одну бесплатную ноду в день}

\hblock{Ограничения} Число подключений --- не нагрузка в мегабитах; связь перегрузки с ошибками из H4 напрямую не установлена (логи только за сентябрь).

\hblock{Вывод} Концентрация бесплатного трафика \textbf{подтверждается}; утверждение, что обвал 27.02 касался только бесплатного пула, \textbf{опровергается} --- это общий сбой сервиса или выгрузки.

\hblock{Рекомендация} Масштабировать бесплатный пул пропорционально числу активных бесплатных пользователей (ориентир --- пользователи на ноду в день), ограничивать скорость на бесплатном тарифе вместо перегрузки ноды, разобрать инцидент 27.02.2026.

%--------------------------------------------------------------
\subsection{H10. Версии приложения и ошибки}

\hblock{Утверждение} Пользователи старых версий приложения чаще сталкиваются с ошибками; принудительное обновление не используется.

\hblock{Данные} Заголовки журнала приложения (платформа, версии ОС и приложения), инциденты из H4, \texttt{mobile\_version}.

\hblock{Метод} Доля инцидентов с недоступностью API и с ошибками VPN-ядра по версиям приложения (не меньше 10 инцидентов на версию), критерий $\chi^2$; разрез по версиям Android.

\hblock{Результат}
\begin{itemize}
  \item Недоступность API: версия 2026.7.16 --- 94,6\,\% инцидентов ($n=56$), 2026.8.19 --- 85,7\,\% ($n=35$), 2026.9.9 --- 74,3\,\% ($n=781$); $\chi^2=10{,}7$, $p=0{,}005$. Ошибки VPN-ядра --- 98,2 / 100 / 86,6\,\%, $p=0{,}001$.
  \item С мая 2026 ни одна версия Android не выпускалась с обязательным обновлением (\texttt{requiredUpdate = false}).
\end{itemize}

\hblock{Ограничения} Все архивы собраны 13--28.09.2026, поэтому старые версии --- это пользователи, не обновившиеся до 2026.9.9; они могут отличаться и устройствами, и сетью. Малые выборки по старым версиям.

\hblock{Вывод} Гипотеза \textbf{слабо подтверждается}: различие статистически значимо, но может объясняться тем, кто не обновляется, а не самой версией.

\hblock{Рекомендация} Включить мягкое, а затем обязательное обновление для версий старше 2026.9.9 и сравнить долю инцидентов до и после.

%==============================================================
\section{Корреляция и причинность}

Все данные наблюдательные, контролируемых экспериментов нет. Поэтому:
\begin{itemize}
  \item для H2 (реклама и покупка), H7 (автопродление), H8 (бывшие гости) и H10 (версии) сделан вывод о \emph{связи}; указаны вероятные смешивающие факторы (вовлечённость, мотивация, самоотбор, время);
  \item там, где возможно, порядок событий задан так, чтобы следствие не влияло на причину: в H2 реклама считается до начала окна покупки, в H3 время покупки сравнивается с концом триала с точностью до минут;
  \item сравнения когорт до и после июня 2026 оговариваются из-за смены состава клиентов (гостевой режим);
  \item для решений по H2, H5 и H7 предложены A/B-тесты.
\end{itemize}

%==============================================================
\section{Заключение}

Наибольший потенциал для продукта, по оценке на этих данных:
\begin{enumerate}
  \item \textbf{Платежи (H5, H6)}: каждая вторая первая оплата не проходит из-за нехватки средств, две трети таких пользователей не возвращаются; июньский сбой продлений стоил около 500 подписчиков сверх обычного оттока.
  \item \textbf{Удержание (H7)}: автопродление связано с продлением в 2,7 раза чаще, но больше половины отключает его в первый месяц.
  \item \textbf{Надёжность подключения (H4, H9, H10)}: типичный инцидент --- недоступность API-домена через DNS, а не отказ нод; бесплатный пул работает с нагрузкой в сотни раз выше платного.
  \item \textbf{Модель монетизации (H2, H3, H8)}: покупают сразу или не покупают вовсе; лимит 500~МБ ограничивает около 30\,\% дней и поддерживает рекламу; триал фактически работает как автосписание.
\end{enumerate}
Гипотеза о массовом мультиаккаунтинге ради триала (H1) не подтвердилась: платный триал уже отсекает злоупотребления.

\end{document}
